\documentclass[12pt,a4paper]{article}

\usepackage[utf8]{inputenc}
\usepackage[T1]{fontenc}
\usepackage[french]{babel}
\usepackage{geometry}
\usepackage{hyperref}
\usepackage{listings}
\usepackage{xcolor}

\geometry{margin=2.5cm}

\definecolor{codegray}{rgb}{0.95,0.95,0.95}
\lstset{
  backgroundcolor=\color{codegray},
  basicstyle=\ttfamily\small,
  breaklines=true,
  frame=single,
  columns=fullflexible
}

\title{FanFlight\\Dossier projet}
\author{Equipe FanFlight}
\date{\today}

\begin{document}

\maketitle
\tableofcontents
\newpage

\section{Presentation du projet}

FanFlight est une application de recherche de vols autour des matchs de la Coupe du Monde 2026.
L'objectif est de proposer une experience simple pour choisir une ville de depart, consulter des matchs disponibles et visualiser des propositions de vols associees.

Le projet regroupe plusieurs briques techniques :

\begin{itemize}
  \item un frontend Next.js pour l'interface utilisateur ;
  \item une API FastAPI pour exposer les donnees de matchs et de vols ;
  \item une base PostgreSQL pour stocker les donnees applicatives ;
  \item un pipeline Spark pour transformer les donnees de vols ;
  \item une chaine CI/CD GitHub Actions vers Docker Hub et une instance EC2.
\end{itemize}

\section{Architecture applicative}

L'application est decoupee en services Docker.

\begin{itemize}
  \item \textbf{frontend} : application Next.js exposee en local sur le port \texttt{3000} et en production sur le port HTTP \texttt{80}.
  \item \textbf{backend-api} : API FastAPI exposee sur le port \texttt{8000}.
  \item \textbf{db} : base PostgreSQL utilisee par l'API et Spark.
  \item \textbf{spark-master} : noeud Spark master, avec listener applicatif.
  \item \textbf{spark-worker} : worker Spark connecte au master.
  \item \textbf{watchtower} : service de production qui met a jour les conteneurs quand une nouvelle image Docker Hub est disponible.
\end{itemize}

L'image PostgreSQL est pinnee sur \texttt{postgres:17}. Il ne faut pas utiliser \texttt{postgres:latest} pour une base de donnees en production, car un changement de version majeure peut rendre le volume existant incompatible sans migration explicite avec \texttt{pg\_upgrade}.

\section{Images Docker}

Les images internes du projet sont publiees sur Docker Hub dans le namespace \texttt{tezzosyris666}.

\begin{lstlisting}
tezzosyris666/fanflight-frontend
tezzosyris666/fanflight-backend-api
tezzosyris666/fanflight-spark
\end{lstlisting}

Chaque image est publiee avec deux tags :

\begin{itemize}
  \item \texttt{latest}, utilise par l'EC2 et Watchtower ;
  \item \texttt{<git-sha>}, utilise pour retrouver une version precise.
\end{itemize}

\section{Configuration locale}

Avant de lancer le projet en local, il faut creer un fichier d'environnement a partir du modele fourni.

\begin{lstlisting}[language=bash]
cp .env.example .env
\end{lstlisting}

Le fichier \texttt{.env.example} contient les variables minimales :

\begin{lstlisting}
DB_NAME=airline_data
DB_USER=root
DB_PASSWORD=password_test
DB_HOST=db
DB_PORT=5432
API_KEY_SERAPI=
\end{lstlisting}

La variable \texttt{API\_KEY\_SERAPI} doit etre renseignee si les appels a l'API externe SerpAPI sont necessaires.

\section{Lancement local}

Le fichier \texttt{docker-compose.local.yml} permet de tester une configuration proche de la production, mais en buildant les images depuis le code local.
Il lance PostgreSQL, Spark master, Spark worker, l'API et le frontend.

\subsection{Demarrer tous les services}

\begin{lstlisting}[language=bash]
docker compose -f docker-compose.local.yml up --build
\end{lstlisting}

\subsection{Verifier les services}

\begin{lstlisting}[language=bash]
curl http://localhost:3000
curl http://localhost:8000/docs
docker compose -f docker-compose.local.yml ps
\end{lstlisting}

Les URLs utiles en local sont :

\begin{itemize}
  \item frontend : \url{http://localhost:3000}
  \item documentation API FastAPI : \url{http://localhost:8000/docs}
  \item Spark master : \url{http://localhost:8080}
  \item Spark worker : \url{http://localhost:8081}
\end{itemize}

\subsection{Consulter les logs}

\begin{lstlisting}[language=bash]
docker compose -f docker-compose.local.yml logs -f db
docker compose -f docker-compose.local.yml logs -f backend-api
docker compose -f docker-compose.local.yml logs -f spark-master
docker compose -f docker-compose.local.yml logs -f spark-worker
\end{lstlisting}

\subsection{Arreter et nettoyer}

Pour arreter les conteneurs sans supprimer les donnees :

\begin{lstlisting}[language=bash]
docker compose -f docker-compose.local.yml down
\end{lstlisting}

Pour supprimer aussi la base PostgreSQL locale :

\begin{lstlisting}[language=bash]
docker compose -f docker-compose.local.yml down -v
\end{lstlisting}

\section{CI/CD}

Le workflow GitHub Actions \texttt{.github/workflows/frontend-deploy.yml} se lance lors d'un push sur \texttt{main} ou \texttt{master} si les dossiers applicatifs sont modifies.

Le pipeline execute les etapes suivantes :

\begin{enumerate}
  \item recuperation du code source ;
  \item connexion a Docker Hub ;
  \item build et push des images frontend, API et Spark ;
  \item connexion SSH a l'instance EC2 ;
  \item generation du fichier \texttt{/opt/fanflight/docker-compose.yml} ;
  \item demarrage ou mise a jour des services avec Docker Compose.
\end{enumerate}

Les secrets GitHub requis sont :

\begin{lstlisting}
DOCKERHUB_USERNAME=tezzosyris666
DOCKERHUB_TOKEN=<token Docker Hub>
RSA=<cle privee SSH de l'EC2>
\end{lstlisting}

\section{Deploiement EC2}

L'instance EC2 cible est accessible avec :

\begin{lstlisting}[language=bash]
ssh -i project.pem ubuntu@35.181.62.34
\end{lstlisting}

En production, les fichiers applicatifs sont stockes dans :

\begin{lstlisting}
/opt/fanflight/docker-compose.yml
/opt/fanflight/.env
\end{lstlisting}

Le fichier \texttt{/opt/fanflight/.env} est conserve entre les deploiements.
Il faut donc modifier ce fichier directement sur l'EC2 pour changer les variables de production.

\section{Verification production}

Apres un deploiement, les commandes utiles sont :

\begin{lstlisting}[language=bash]
ssh -i project.pem ubuntu@35.181.62.34
cd /opt/fanflight
sudo docker compose ps
sudo docker logs fanflight_frontend
sudo docker logs backend_api
sudo docker logs spark-master
sudo docker logs spark-worker
sudo docker logs fanflight_watchtower
\end{lstlisting}

Le frontend est expose sur :

\begin{lstlisting}
http://35.181.62.34
\end{lstlisting}

Le security group AWS doit autoriser le port TCP \texttt{80} depuis Internet.
Les autres ports exposes par Docker Compose ne doivent etre ouverts que si un besoin technique le justifie.

\section{Points d'attention}

\begin{itemize}
  \item Les repositories Docker Hub doivent exister avant le premier push d'image.
  \item Le secret GitHub \texttt{RSA} doit contenir la cle privee complete, sans guillemets.
  \item Le test local doit etre lance avant un push en production.
  \item En production, Watchtower met a jour les conteneurs qui suivent le tag \texttt{latest}.
  \item Le volume PostgreSQL local peut etre supprime avec \texttt{docker compose -f docker-compose.local.yml down -v}.
\end{itemize}

\end{document}
